\section{Quasi-affine schemes}
\label{section-quasi-affine}

\begin{definition}
\label{definition-quasi-affine}
A scheme $X$ is called {\it quasi-affine} if it is quasi-compact
and isomorphic to an open subscheme of an affine scheme.
\end{definition}

\begin{lemma}
\label{lemma-quasi-coherent-quasi-affine}
Let $A$ be a ring and let $U \subset \Spec(A)$ be a quasi-compact
open subscheme. For $\mathcal{F}$ quasi-coherent on $U$ the canonical map
$$
\widetilde{H^0(U, \mathcal{F})}|_U \to \mathcal{F}
$$
is an isomorphism.
\end{lemma}

\begin{proof}
Denote $j : U \to \Spec(A)$ the inclusion morphism. Then
$H^0(U, \mathcal{F}) = H^0(\Spec(A), j_*\mathcal{F})$ and
$j_*\mathcal{F}$ is quasi-coherent by
Schemes, Lemma \ref{schemes-lemma-push-forward-quasi-coherent}.
Hence $j_*\mathcal{F} = \widetilde{H^0(U, \mathcal{F})}$ by
Schemes, Lemma \ref{schemes-lemma-equivalence-quasi-coherent}.
Restricting back to $U$ we get the lemma.
\end{proof}

\begin{lemma}
\label{lemma-invert-f-affine}
Let $X$ be a scheme. Let $f \in \Gamma(X, \mathcal{O}_X)$.
Assume $X$ is quasi-compact and quasi-separated and assume that
$X_f$ is affine. Then the canonical morphism
$$
j : X \longrightarrow \Spec(\Gamma(X, \mathcal{O}_X))
$$
from Schemes, Lemma \ref{schemes-lemma-morphism-into-affine}
induces an isomorphism of $X_f = j^{-1}(D(f))$ onto the standard affine
open $D(f) \subset \Spec(\Gamma(X, \mathcal{O}_X))$.
\end{lemma}

\begin{proof}
This is clear as $j$ induces an isomorphism of rings
$\Gamma(X, \mathcal{O}_X)_f \to \mathcal{O}_X(X_f)$ by
Lemma \ref{lemma-invert-f-sections} above.
\end{proof}

\begin{lemma}
\label{lemma-quasi-affine}
Let $X$ be a scheme. Then $X$ is quasi-affine if and only if
the canonical morphism
$$
X \longrightarrow \Spec(\Gamma(X, \mathcal{O}_X))
$$
from Schemes, Lemma \ref{schemes-lemma-morphism-into-affine} is
a quasi-compact open immersion.
\end{lemma}

\begin{proof}
If the displayed morphism is a quasi-compact open immersion then
$X$ is isomorphic to a quasi-compact open subscheme of
$\Spec(\Gamma(X, \mathcal{O}_X))$ and clearly $X$ is quasi-affine.

\medskip\noindent
Assume $X$ is quasi-affine, say $X \subset \Spec(R)$ is
quasi-compact open. This in particular implies that $X$ is
separated, see
Schemes, Lemma \ref{schemes-lemma-subscheme-of-separated-scheme}.
Let $A = \Gamma(X, \mathcal{O}_X)$.
Consider the ring map $R \to A$ coming from
$R = \Gamma(\Spec(R), \mathcal{O}_{\Spec(R)})$
and the restriction mapping of the sheaf $\mathcal{O}_{\Spec(R)}$.
By Schemes, Lemma \ref{schemes-lemma-morphism-into-affine}
we obtain a factorization:
$$
X \longrightarrow
\Spec(A) \longrightarrow
\Spec(R)
$$
of the inclusion morphism. Let $x \in X$. Choose $r \in R$ such that
$x \in D(r)$ and $D(r) \subset X$. Denote $f \in A$ the image of $r$
in $A$. The open $X_f$ of Lemma \ref{lemma-invert-f-sections}
above is equal to $D(r) \subset X$ and hence $A_f \cong R_r$ by the
conclusion of that lemma.
Hence $D(r) \to \Spec(A)$ is an isomorphism onto the
standard affine open $D(f)$ of $\Spec(A)$. Since $X$
can be covered by such affine opens $D(f)$ we win.
\end{proof}

\begin{lemma}
\label{lemma-cartesian-diagram-quasi-affine}
Let $U \to V$ be an open immersion of quasi-affine schemes. Then
$$
\xymatrix{
U \ar[d] \ar[rr]_-j & & \Spec(\Gamma(U, \mathcal{O}_U)) \ar[d] \\
U \ar[r] & V \ar[r]^-{j'} & \Spec(\Gamma(V, \mathcal{O}_V))
}
$$
is cartesian.
\end{lemma}

\begin{proof}
The diagram is commutative by Schemes, Lemma
\ref{schemes-lemma-morphism-into-affine}.
Write $A = \Gamma(U, \mathcal{O}_U)$ and $B = \Gamma(V, \mathcal{O}_V)$. Let
$g \in B$ be such that $V_g$ is affine and contained in $U$. This
means that if $f$ is the image of $g$ in $A$, then $U_f = V_g$. By Lemma
\ref{lemma-invert-f-affine} we see that $j'$ induces an isomorphism of
$V_g$ with the standard open $D(g)$ of $\Spec(B)$. Thus
$V_g \times_{\Spec(B)} \Spec(A) \to \Spec(A)$ is an
isomorphism onto $D(f) \subset \Spec(A)$. By Lemma \ref{lemma-invert-f-affine}
again $j$ maps $U_f$ isomorphically to $D(f)$. Thus we see that
$U_f = U_f \times_{\Spec(B)} \Spec(A)$. Since by
Lemma \ref{lemma-quasi-affine} we can cover $U$ by $V_g = U_f$ as above,
we see that $U \to U \times_{\Spec(B)} \Spec(A)$ is an isomorphism.
\end{proof}

\begin{lemma}
\label{lemma-quasi-affine-presentation}
Let $X$ be a quasi-affine scheme. There exists an integer $n \geq 0$,
an affine scheme $T$, and a morphism $T \to X$ such that for every
morphism $X' \to X$ with $X'$ affine the fibre product $X' \times_X T$
is isomorphic to $\mathbf{A}^n_{X'}$ over $X'$.
\end{lemma}

\begin{proof}
By definition, there exists a ring $A$ such that $X$ is isomorphic to a
quasi-compact open subscheme $U \subset \Spec(A)$. Recall that the standard
opens $D(f) \subset \Spec(A)$ form a basis for the topology, see
Algebra, Section \ref{algebra-section-spectrum-ring}. Since $U$ is
quasi-compact we can choose $f_1, \ldots, f_n \in A$ such that
$U = D(f_1) \cup \ldots \cup D(f_n)$. Thus we may assume
$X = \Spec(A) \setminus V(I)$ where $I = (f_1, \ldots, f_n)$. We set
$$
T = \Spec(A[t, x_1, \ldots, x_n]/(f_1 x_1 + \ldots + f_n x_n - 1))
$$
The structure morphism $T \to \Spec(A)$ factors through the open $X$
to give the morphism $T \to X$. If $X' = \Spec(A')$ and the morphism
$X' \to X$ corresponds to the ring map $A \to A'$, then the images
$f'_1, \ldots, f'_n \in A'$ of $f_1, \ldots, f_n$
generate the unit ideal in $A'$.
Say $1 = f'_1 a'_1 + \ldots + f'_n a'_n$.
The base change $X' \times_X T$ is the spectrum of
$A'[t, x_1, \ldots, x_n]/(f'_1 x_1 + \ldots + f'_n x_n - 1)$.
We claim the $A'$-algebra homomorphism
$$
\varphi :
A'[y_1, \ldots, y_n]
\longrightarrow
A'[t, x_1, \ldots, x_n, x_{n + 1}]/(f'_1 x_1 + \ldots + f'_n x_n - 1)
$$
sending $y_i$ to $a'_i t + x_i$ is an isomorphism. The claim finishes
the proof of the lemma. The inverse of $\varphi$ is given by the $A'$-algebra
homomorphism
$$
\psi :
A'[t, x_1, \ldots, x_n, x_{n + 1}]/(f'_1 x_1 + \ldots + f'_n x_n - 1)
\longrightarrow
A'[y_1, \ldots, y_n]
$$
sending $t$ to $-1 + f'_1 y_1 + \ldots + f'_n y_n$ and $x_i$ to
$y_i + a'_i - a'_i(f'_1 y_1 + \ldots + f'_n y_n)$ for $i = 1, \ldots, n$.
This makes sense because $\sum f'_ix_i$ is mapped to
$$
\begin{matrix}
\sum f'_i(y_i + a'_i - a'_i(\sum f'_j y_j)) =
(\sum f'_iy_i) + 1 - (\sum f'_j y_j) = 1
\end{matrix}
$$
To see the maps are mutually inverse one computes as follows:
$$
\begin{matrix}
\varphi(\psi(t) = \varphi(-1 +  \sum f'_i y_i) =
-1 + \sum f'_i (a'_i t + x_i) = t \\
\varphi(\psi(x_i)) = \varphi(y_i + a'_i - a'_i(\sum f'_j y_j)) =
a'_i t + x_i + a'_i - a'_i(\sum f'_ja'_jt + f'_jx_j) = x_i \\
\psi(\varphi(y_i)) = \psi(a'_i t + x_i) =
a'_i(-1 + \sum f'_j y_j) + y_i + a'_i - a'_i(\sum f'_j y_j) = y_i
\end{matrix}
$$
This finishes the proof.
\end{proof}




